\addchap{\lsAcknowledgementTitle}

Ich danke Hubert Haider, Matthias Hüning und Shravan Vasishth für die Diskussion sowie Martin Haspelmath und Tibor Kiss für ausführliche Kommentare
zu einer früheren Version dieses Buches. Ich danke Antonio Machicao y Priemer für Kommentare zu den Abschnitten
über alternative Ansätze. Ich danke Shravan Vasishth für die Diskussion psycholinguistischer Fragen.

Ich möchte allen Studierenden meines Kurses über die Struktur der germanischen Sprachen danken. Dieses Buch
hat enorm von ihren Fragen und der Diskussion während der Vorlesungen profitiert.
Carolin Ulmer %05.09. WALS falsch zitiert.
und
Bastian Schoenfeld % 14.06.2020 zur Behauptung, dass die Argumente im Englischen alle vollständig
                   % sein müssen, was für VPen nicht stimmt.
verdienen für wichtige Kommentare besondere Erwähnung. Ich danke
Sarah Böke,
Alexandra Fosså,
Anne Kilgus,
Nils P. Kujath
und
Maya Tabea Wolff
für Hinweise auf Tippfehler. Ich danke Hannah Schleupner für die Hilfe bei Abbildung~\ref{fig-history-relations-germanic}.

Ich danke der Abteilung für Englische Linguistik an der Universität Frankfurt/""Main für die Einladung, über
den diesem Buch zugrunde liegenden Ansatz zur germanischen Syntax zu sprechen.

Ich möchte allen Korrekturleserinnen und Korrekturlesern danken (\makeatletter\@proofreader\makeatother), die weit mehr als
übliches Lektorat geleistet und auch zum Inhalt Rückmeldungen gegeben haben. Besonderer Dank geht an Brett Reynolds und Anne Zaenen.

Dieses Buch wurde teilweise von der Deutschen Forschungsgemeinschaft (DFG) gefördert – SFB 1412, 416591334, Projekt A04.

Ich danke der Twitter-Community für die Diskussion über Geschlechterfragen und Namen in linguistischen Beispielen mit
mir. Besonderer Dank geht an Lauren Ackerman, Kirby Conrod und Brianna Wilson für die ausführliche Diskussion.


Elisabeth Eberle und Sigi Bahr haben das Buch im Tutorium unserer
Vorlesung zur germanischen Syntax eingesetzt. Danke für Kommentare, Fragen und die Diskussion.

Eine Rohfassung der deutschen Übersetzung dieses Buches wurde mit künstlicher Intelligenz (Claude AI) erstellt. Ich danke meinem Bruder Frank
J.\ Müller für die Hilfe dabei. Claude AI hat den Haupttext und die Glossierungen in den Beispielen
und Bäumen automatisch übersetzt und die Glossierungen und Übersetzungen der deutschen Beispiele
entfernt. Ich hätte nie geglaubt, dass ich dieses Buch einmal ins Deutsche übertragen würde, aber
die gemeinsame Arbeit mit Frank an Millions4earth\footnote{
\url{https://millions4.earth/m/ATZM4X727Y?lang=de}, 24.06.2026.
} hat mir gezeigt, was mit KI inzwischen möglich ist. Der ursprüngliche Prompt war \emph{Kannst du das übersetzen, ohne es kaputt zu machen?}. Die
gesamte Interaktion ist hier dokumentiert: \url{https://github.com/langsci/Syntax-der-germanischen-Sprachen/blob/main/Comments/chat-protokoll-2026-06-24.md}.



%      <!-- Local IspellDict: de_DE -->
